% ============================================================
%  ShARC — Architecture Novelty & Literature Comparison
%  Draft v1.0  |  Agriya  |  April 2026
% ============================================================
%
%  Compile with: pdflatex sharc_arch_novelty.tex
% ============================================================

\documentclass[11pt]{article}

\usepackage{amsmath, amssymb}
\usepackage{booktabs}
\usepackage{geometry}
\usepackage{xcolor}
\usepackage{hyperref}

\geometry{margin=1.1in}

\newcommand{\CVaR}{\operatorname{CVaR}}
\newcommand{\todo}[1]{\textcolor{red}{\textbf{[TODO: #1]}}}

\title{\textbf{ShARC: Architecture Novelty \& Literature Comparison}\\
\large How ShARC Differs from Existing Neural Combinatorial Optimisation Architectures}
\author{Agriya \\ \small Ashoka University}
\date{Draft --- \today}

\begin{document}
\maketitle

% ============================================================
\section{Overview}
% ============================================================

Neural combinatorial optimisation (NCO) for routing has converged on a
standard encoder-decoder transformer architecture since Kool et
al.~\cite{kool2019attention}.  ShARC adopts this backbone but departs from
it in three principled ways that together address a research gap no existing
work has closed: \emph{risk-sensitive arc routing under compound distribution
shifts}.  This document situates ShARC within the NCO architecture landscape
and identifies exactly where and why it is novel.

% ============================================================
\section{The Standard NCO Architecture and Its Limitations}
% ============================================================

\paragraph{Kool et al. (ICLR 2019)~\cite{kool2019attention}.}
The attention model for CVRP uses a transformer encoder over node coordinates
and a masked cross-attention pointer decoder.  It is trained via REINFORCE
with a greedy rollout baseline.  This architecture is the foundation for
virtually all subsequent NCO work and is the direct ancestor of ShARC's
policy backbone.

\textbf{Limitations for ShARC's setting:} (i)~it assumes a fixed,
stationary instance distribution — the encoder is never conditioned on any
shift signal; (ii)~it optimises expected tour length, not tail risk;
(iii)~it is designed for node routing (VRP/TSP), not arc routing (CARP),
where actions correspond to edges rather than nodes and feasibility masking
must account for service demand along arcs.

% ============================================================
\section{Distribution-Shift NCO Architectures}
% ============================================================

\paragraph{DRO-VRP (AAAI 2022)~\cite{jiang2022dro}.}
Applies group distributionally robust optimisation (group DRO) to the
REINFORCE loss.  The model is a standard Kool et al.\ attention model;
the DRO modification is purely at the loss level, reweighting training
instances to focus on worst-case distribution groups.

\textbf{Gap vs. ShARC:} DRO-VRP treats distribution shift as movement
between a discrete set of pre-defined instance families.  It does not
condition the policy on the current shift, does not model a continuous
multi-parameter shift vector, and does not optimise a tail risk measure.
The uncertainty set is over node coordinate distributions (rectangular),
not over demand, cost, and availability jointly (non-rectangular).

\paragraph{Omni-VRP (ICML 2023)~\cite{zhou2023omni}.}
A meta-learning framework for generalisation across both instance size and
distribution.  Uses MAML-style inner/outer loops to adapt the attention
model to new distributions at test time.

\textbf{Gap vs. ShARC:} Omni-VRP addresses cross-distribution
generalisation but does not provide tail-risk guarantees and does not
condition the policy on a shift context vector.  The uncertainty is over
discrete distribution families, not a continuous shift parameterisation.

\paragraph{SHIELD (ICML 2025)~\cite{goh2025shield}.}
A multi-task, multi-distribution VRP solver using sparsity and hierarchical
decoding to generalise across discretely enumerated VRP variants and
distribution families simultaneously.

\textbf{Gap vs. ShARC:} SHIELD handles a fixed, discrete set of
distribution shifts and does not model continuous drift.  Its hierarchical
decoding is designed for multi-task generalisation (different constraint
types), not for tail-risk-sensitive behaviour under continuously shifting
parameters.  SHIELD has no CVaR objective or risk-sensitive training signal.

\paragraph{Collaboration (NeurIPS 2024)~\cite{zhou2024collaboration}.}
An ensemble-based adversarial training framework that improves robustness
of NCO models to input perturbations.  Multiple models are trained
collaboratively to defend against crafted adversarial attacks on instance
features.

\textbf{Gap vs. ShARC:} Collaboration addresses adversarial robustness
(worst-case input perturbations by an adversary) rather than distributional
robustness (continuous stochastic shifts).  It does not model arc
availability dropout, does not use a risk measure, and is designed for
node routing.

% ============================================================
\section{Risk-Sensitive RL Architectures}
% ============================================================

\paragraph{Tamar et al. (AAAI 2015)~\cite{tamar2015optimizing}.}
Derives the CVaR-REINFORCE gradient estimator for general MDPs and
demonstrates it on small tabular and continuous control tasks.  No
encoder-decoder architecture, no routing domain, no shift conditioning.

\paragraph{Chow et al. (JMLR 2018)~\cite{chow2018risk}.}
Actor-critic algorithms for CVaR-constrained MDPs via Lagrangian relaxation.
The constraint is on the CVaR of cumulative cost; satisfaction is guaranteed
in expectation but not in the tail under distribution shift.

\textbf{Gap vs. ShARC:} Neither Tamar nor Chow et al.\ combine risk-sensitive
RL with a transformer-based combinatorial optimisation policy or with a
distribution shift model.  The connection between CVaR objectives and routing
generalisation under compound shifts has not been studied.

% ============================================================
\section{Arc Routing Neural Methods}
% ============================================================

No prior NCO work addresses arc routing specifically.  The closest is the
UCARP literature~\cite{mei2010carp}, which uses genetic programming
hyper-heuristics to evolve routing policies under stochastic arc demands.

\textbf{Gap vs. ShARC:} GP-based UCARP methods (i)~assume a fixed stochastic
environment (demands drawn from fixed per-arc distributions), not a
continuously drifting shift; (ii)~cannot be conditioned on a shift context
vector; (iii)~do not optimise a tail risk measure; and (iv)~are not
differentiable, precluding gradient-based training.

% ============================================================
\section{ShARC's Architectural Contributions}
% ============================================================

Table~\ref{tab:arch-comparison} summarises the architectural differences
between ShARC and the most closely related prior work.

\begin{table}[h!]
\centering
\small
\caption{Architectural comparison of ShARC against related NCO and risk-sensitive RL methods.
\checkmark = present; -- = absent.}
\label{tab:arch-comparison}
\begin{tabular}{lccccc}
\toprule
\textbf{Method} & \textbf{Arc routing} & \textbf{Shift context} & \textbf{CVaR objective} &
\textbf{Continuous shift} & \textbf{Budget signal} \\
\midrule
Kool et al.~\cite{kool2019attention}    & -- & -- & -- & -- & -- \\
DRO-VRP~\cite{jiang2022dro}             & -- & -- & -- & -- & -- \\
Omni-VRP~\cite{zhou2023omni}            & -- & -- & -- & -- & -- \\
SHIELD~\cite{goh2025shield}             & -- & -- & -- & -- & -- \\
Collaboration~\cite{zhou2024collaboration} & -- & -- & -- & -- & -- \\
Tamar et al.~\cite{tamar2015optimizing} & -- & -- & \checkmark & -- & -- \\
Chow et al.~\cite{chow2018risk}         & -- & -- & \checkmark & -- & -- \\
UCARP / GP~\cite{mei2010carp}           & \checkmark & -- & -- & -- & -- \\
\midrule
\textbf{ShARC (ours)} & \checkmark & \checkmark & \checkmark & \checkmark & \checkmark \\
\bottomrule
\end{tabular}
\end{table}

ShARC introduces three specific architectural novelties not present in any
prior work:

\begin{enumerate}
    \item \textbf{Shift context injection into the decoder.}
    The decoder context vector explicitly includes a learned projection of
    the current episode's shift vector $\phivec = (\delta_d, \delta_c,
    \delta_p)$.  This allows the policy to adapt its routing decisions
    to the \emph{type and severity} of the current shift, not just to
    average over shifts.  No prior NCO architecture conditions the decoder
    on a continuous distributional shift signal.

    \item \textbf{Budget variable in the decoder context.}
    The running maximum vehicle time $\beta_t = \max_k \tau_{k,t}$ is
    included in the decoder context at every step.  This gives the policy
    an explicit signal about its current worst-case cost trajectory,
    enabling risk-sensitive routing decisions within the episode --- for
    example, preferring to service high-priority arcs early when the budget
    is already tight.  No prior routing NCO method includes a within-episode
    risk budget signal.

    \item \textbf{CVaR-REINFORCE with tail-normalised gradients for
    combinatorial optimisation.}
    The CVaR-REINFORCE estimator of Tamar et al.\ is adapted to the HCARP
    setting with two modifications: (i)~the greedy rollout CVaR (rather than
    mean) is used as the baseline to ensure tail-consistent advantage
    estimation; and (ii)~advantage normalisation is applied within the
    worst-$\alpha$ tail to prevent gradient magnitude collapse under the
    large raw makespans typical of arc routing instances.  This combination
    has not been applied in any prior NCO work.
\end{enumerate}

\begin{remark}[What ShARC does \emph{not} claim as novel]
The transformer encoder-decoder backbone, pointer network, and REINFORCE
training loop are all inherited from~\cite{kool2019attention} and are not
claimed as contributions.  The novelty is in the specific combination of
(a)~arc routing domain, (b)~continuous compound shift model, (c)~shift
context injection, (d)~budget variable, and (e)~CVaR tail objective.
\end{remark}

% ============================================================
\section{Research Gap Summary}
% ============================================================

The gap ShARC addresses can be stated precisely: \emph{no existing work
optimises a tail risk measure for arc routing under a continuous,
multi-parameter distribution shift, nor conditions a neural routing policy
on an explicit shift context to enable shift-aware decision-making.}

More specifically:
\begin{itemize}
    \item Existing NCO distribution-shift methods (DRO-VRP, SHIELD,
    Collaboration, Omni-VRP) address VRP not CARP, use discrete not
    continuous shift models, and optimise expected not tail performance.
    \item Existing CVaR RL methods (Tamar, Chow) do not apply to
    combinatorial optimisation and do not model distribution shift.
    \item Existing UCARP methods use evolutionary not gradient-based methods,
    assume fixed not drifting environments, and do not condition on shift.
\end{itemize}

ShARC is the first work to sit at all three intersections simultaneously:
arc routing $\cap$ neural RL policy $\cap$ CVaR under continuous distribution
shift.

% ============================================================
%  Bibliography
% ============================================================
\bibliographystyle{plain}
\begin{thebibliography}{9}

\bibitem{kool2019attention}
W.~Kool, H.~van Hoof, and M.~Welling.
\newblock Attention, learn to solve routing problems!
\newblock In \emph{Proc. ICLR}, 2019.

\bibitem{jiang2022dro}
Y.~Jiang, Y.~Wu, Z.~Cao, and J.~Zhang.
\newblock Learning to solve routing problems via distributionally robust
  optimization.
\newblock In \emph{Proc. AAAI}, 2022.

\bibitem{zhou2023omni}
J.~Zhou, Y.~Wu, W.~Song, Z.~Cao, and J.~Zhang.
\newblock Towards omni-generalizable neural methods for vehicle routing
  problems.
\newblock In \emph{Proc. ICML}, 2023.

\bibitem{goh2025shield}
Y.~L.~Goh, Z.~Cao, Y.~Ma, J.~Zhou, M.~H.~Dupty, and W.~S.~Lee.
\newblock {SHIELD}: Multi-task multi-distribution vehicle routing solver with
  sparsity and hierarchy.
\newblock In \emph{Proc. ICML}, 2025.

\bibitem{zhou2024collaboration}
J.~Zhou, Y.~Wu, Z.~Cao, W.~Song, J.~Zhang, and Z.~Shen.
\newblock Collaboration! towards robust neural methods for routing problems.
\newblock In \emph{Proc. NeurIPS}, 2024.

\bibitem{tamar2015optimizing}
A.~Tamar, Y.~Glassner, and S.~Mannor.
\newblock Optimizing the {CVaR} via sampling.
\newblock In \emph{Proc. AAAI}, 2015.

\bibitem{chow2018risk}
Y.~Chow, M.~Ghavamzadeh, L.~Janson, and M.~Pavone.
\newblock Risk-constrained reinforcement learning with percentile risk
  criteria.
\newblock \emph{Journal of Machine Learning Research}, 18(167):1--51, 2018.

\bibitem{mei2010carp}
Y.~Mei, K.~Tang, and X.~Yao.
\newblock Capacitated arc routing problem in uncertain environments.
\newblock In \emph{Proc. IEEE CEC}, 2010.

\bibitem{bello2016neural}
I.~Bello, H.~Pham, Q.~V.~Le, M.~Norouzi, and S.~Bengio.
\newblock Neural combinatorial optimization with reinforcement learning.
\newblock In \emph{Proc. ICLR Workshop}, 2017.

\end{thebibliography}

\end{document}